\FloatBarrier
\section{Three official submissions and what they measure}
\label{sec:official}

\subsection{Private-test results}

Table~\ref{tab:official-submissions} places the three PHASE-TTC submissions beside the documented Garl-TTC reference. Lower MiD is better. The original event--RGB system obtained 35.80 Overall MiD on 28 October 2026. A pre-registered RGB-plus-box-history submission obtained 23.47 on 18 November after removing event inputs. A third, pre-registered comparator on 21 November removed the RGB estimates as well and read only 600\,ms of supplied-box history; it obtained 27.98. All three scoring archives contain the same 6,762 tokens and zero scorer-defined failures. The downloaded third record, including its four band scores, is preserved in repository commit \texttt{8fa8dda6}. Because the public leaderboard retains one row per participant, dated submission records rather than a current leaderboard row preserve the three results.

\begin{table}[htbp]
\centering\footnotesize
\begin{tabularx}{\linewidth}{@{}p{.17\linewidth}p{.12\linewidth}p{.17\linewidth}Xp{.085\linewidth}p{.07\linewidth}@{}}
\toprule
Entry and date & Submission & Sensor inputs & Supplied-region and temporal support & Overall MiD & Failure \\
\midrule
Garl-TTC reference\\28 Jun 2026 & 817519 & Events + RGB & Supplied regions; two observations 0.1 s apart & 45.03 & 0\% \\
PHASE-TTC original\\28 Oct 2026 & 1016512 & Events + RGB & Supplied regions; up to 650 ms causal history & 35.80 & 0\% \\
PHASE-TTC-Geo\\18 Nov 2026 & 1058935 & RGB; no event input & Supplied regions; 600 ms box history within the 650 ms budget & 23.47 & 0\% \\
\bottomrule
\end{tabularx}

\caption{Official GarlTTC Test results over 6,762 tokens. ``Failure'' is the scorer-defined invalid-output rate, not a safety failure rate. The reference and PHASE systems use different models and temporal support.}
\label{tab:official-submissions}
\end{table}

The 23.47 result is 34\% below the original system and 48\% below the reference. The matched box-only comparator is 4.50325 MiD points higher than 23.47238; equivalently, adding the RGB estimates reduces Overall MiD by 16.10\% relative to the 27.97563 comparator. These official rows establish that supplied-region history alone is strong but does not reproduce the RGB-plus-geometry score under the tested recipes. They do not isolate a sensor-level causal effect or establish that supplied boxes replace event measurements under delayed or noisy detections, frame-filling objects, or another acquisition protocol.

\subsection{Metric and band asymmetry}

The scorer maps TTC to phase and reports
\begin{equation}
 \mathrm{MiD}(\widehat T,T)=10^4\left|\log\left(1-\frac{0.1}{T}\right)
 -\log\left(1-\frac{0.1}{\widehat T}\right)\right|.
\end{equation}
It then weights the crucial, small, large, and negative TTC bands by 0.5, 0.3, 0.1, and 0.1. The metric therefore emphasizes approaching contacts below 6\,s and is more sensitive near its positive-domain boundary. It is not an error in seconds. The second submission reduces the crucial-band score from 43.52 to 24.37; that 19.15-point change contributes 9.58 of the 12.33 Overall points. Changes in the small, large, and negative bands contribute 1.29, 0.28, and 1.18 points. The weighted result should not be summarized as a uniform improvement across TTC regimes.

\begin{table}[htbp]
\centering\small
\begin{tabular}{lrrrrr}
\toprule
Entry & Overall & Crucial & Small & Large & Negative \\
\midrule
Garl-TTC reference & 45.03 & 53.13 & 37.58 & 40.61 & 31.29 \\
PHASE-TTC original & 35.80 & 43.52 & 28.26 & 25.72 & 29.93 \\
PHASE-TTC-Geo & \textbf{23.47} & 24.37 & 23.95 & 22.89 & 18.13 \\
Box-history comparator & 27.98 & 26.39 & 28.00 & 32.94 & 30.89 \\
\bottomrule
\end{tabular}
\caption{Official band scores. Overall uses the 0.5/0.3/0.1/0.1 band weights rather than a sample-weighted mean.}
\label{tab:official-submission-bands}
\end{table}

The 6,762 tokens represent 5,805 physical queries and 957 aliases. Aliases reuse predictions at the same sensor query. The scorer flags nonfinite TTC or $|\widehat T|<0.1$\,s as failures, then averages finite MiD values by band. Zero scorer-defined failures therefore certifies finite admissible outputs under that rule; it does not certify calibrated risk, localization, or safe operation.

\subsection{Comparator and modality boundaries}

The official reference row comes from the leaderboard and was not reproduced locally for Table~\ref{tab:official-submissions}. Garl-TTC and the original PHASE-TTC both use events, RGB, and supplied regions, but Garl-TTC estimates a height ratio from two observations 0.1\,s apart while PHASE-TTC aggregates up to 650\,ms of branch history. The benchmark does not standardize temporal support, training data, learned components, or runtime. The 35.80-versus-45.03 difference belongs to the complete systems.

Separate local controls use the released Garl checkpoints under their native input contract. They select the latest causal RGB endpoint and the preceding endpoint nearest 100\,ms, construct a polarity-agnostic 20-plane event time volume at each endpoint with at most 100\,ms support, and use the supplied regions to form a common square crop. The adapter substitutes the measured endpoint separation in the height-ratio-to-TTC conversion; it does not refit the checkpoint. Full, event-only, and visual-only routes have different input availability, and these development-exposed runs are not the official leaderboard comparison (Appendix~\ref{app:competition-details}). In the later service experiment, a CPU path constructs Garl event volumes on cache misses and hashes prepared inputs inside the measured call. Its timing is specific to that implementation and is not a lower bound for Garl-TTC.

The geometry submission changes both modality and measurement access. It uses no event input, but it receives 600\,ms of causal box expansion in addition to RGB-branch estimates. The leaderboard fact sheet displays its ``Uses Event'' field as empty rather than false. We retain that downloaded record and derive the no-event statement from the bound recipe and scored package, not from the empty form field. The third submission uses the same supplied-box history and a learned geometry head but no RGB image estimates and no event input. It is an official matched comparator, not a replacement release for PHASE-TTC-Geo.

\subsection{Selection, exposure, and attribution}

The original system was the author's first test submission. The geometry variant was chosen by a rule fixed on TRAIN40 development folds and then submitted once with a receipt that bound its code, seed, package, and hypothesis. Those folds showed 85.17 Overall MiD for the released recipe and 41.27 for the selected geometry head, but fold and server levels differ enough that these values must not be pooled or used as a calibration curve. The private result confirms the direction of the geometry intervention, not the size predicted by the folds. The box-only comparator was likewise bound before its third upload and was reported whether or not it beat 23.47. Its 27.98 score provides a private-test ablation of RGB estimates conditional on the fitted heads and shared box-history contract.

Only the three private-test scores use targets that the project did not inspect during development. They are sequential exposures to the same server, not independent test sets: the earlier score was known before each later recipe was uploaded. Transfer, routing, event-history, literature-baseline, and FCWD controls reuse inspected recordings or open labels. They identify failure modes and plausible mechanisms without providing another independent test. Appendix~\ref{sec:post-audit-geometry} records the second submission and its fold analysis; the third submission record and pre-registered matched-comparator protocol preserve the box-only comparison. Appendix~\ref{sec:followup} records later controls and their exposure.

\subsection{Original submission by band}

Table~\ref{tab:official-benchmark} retains the full-precision calculation for the original submission and the dated reference. Percentage reductions combine the unrounded PHASE-TTC scorer values with the reference's displayed two-decimal values.

\begin{table}[H]
\centering\small
\caption{Official GarlTTC Test scores on 28 October 2026 (lower is better). PHASE-TTC scores are shown to four decimals and Garl-TTC at published precision. Relative reductions use unrounded PHASE-TTC scores. Both entries have 6,762 scored tokens and 0\% scorer-defined failures. \evidence{B01,B02}}
\label{tab:official-benchmark}
\begin{tabular}{lrrr}
\toprule
\textbf{Metric} & \textbf{PHASE-TTC} & \textbf{Garl-TTC} & \textbf{Decrease} \\
\midrule
Overall & 35.8015 & 45.03 & 20.49\% \\
Crucial & 43.5190 & 53.13 & 18.09\% \\
Small & 28.2568 & 37.58 & 24.81\% \\
Large & 25.7179 & 40.61 & 36.67\% \\
Negative & 29.9314 & 31.29 & 4.34\% \\
\bottomrule
\end{tabular}
\end{table}


For band $b$, the additive contribution to the original-versus-reference difference is $w_b(\mathrm{MiD}^{\mathrm{Garl}}_b-\mathrm{MiD}^{\mathrm{PHASE}}_b)$. The crucial and small bands contribute \OfficialCrucialSmallContributionShare{} of that Overall difference. This decomposition reflects the benchmark weights and must remain separate from sample-weighted mean MiD or TTC error in seconds.

The current platform does not regenerate the historical original predictions byte for byte. A native replay changes 504 of 5,805 physical-query TTC values beyond the historical tolerance. Converting the largest admissible displacement of each replayed output to the benchmark coordinate gives a mean bound of 0.010 MiD and a maximum of 0.30 for one query, without access to private targets. The scored package and its official result remain the evidence for 35.80; local repeatability does not replace them.
\FloatBarrier
